\section{Discussion}
\label{sec:discussion}

The comparative analysis of the machine learning classifiers provides insight not only into performance differences, but also into how model behavior interacts 
with the low-dimensional, subject-level feature space used in this study. Because the proposed framework relies on aggregated extremal gyroscopic features 
rather than high-frequency temporal representations, model performance is primarily governed by robustness to inter-subject variability, sensitivity to noise, 
and the ability to generalize from limited samples.

Among all models, \textit{Random Forest}, \textit{XGBoost}, and \textit{Logistic Regression} emerged as the top-performing classifiers across nearly all evaluation 
metrics. The \textit{Random Forest} model achieved the highest F1-score (93\%) and near-perfect classification in the confusion matrix, with only one instance 
misclassified. This model benefited from its ensemble learning mechanism, which combines multiple decision trees to reduce overfitting and capture complex, 
non-linear patterns in the data. Similarly, \textit{XGBoost} showed excellent generalization, with a high AUC of 0.96 and minimal misclassification, thanks to 
its gradient-boosted tree approach that optimizes classification boundaries based on prior errors. \textit{Logistic Regression}, while a linear model, also 
performed competitively (F1-score of 89\%), likely due to the high discriminative power of the selected features and the low-dimensional nature of the dataset.

In contrast, the \textit{SVMs}, especially the polynomial kernel variant, underperformed in comparison. While SVMs 
with linear and RBF kernels achieved reasonable AUC values (0.92 and 1.00 respectively), their F1-scores remained moderate (85\%), and they showed a higher 
false negative rate. This tendency to misclassify abnormal gait as normal is especially problematic in clinical applications where sensitivity to pathological cases 
is critical. The observed performance gap can be attributed to the SVMs’ sensitivity to kernel parameterization and decision boundary rigidity when operating on small, 
low-dimensional datasets, where minor variations in feature distributions can lead to suboptimal margin placement. Lastly, the \textit{KNN} model displayed the 
lowest overall performance. Its reliance on instance-based learning likely made it more susceptible to local noise and feature overlap between classes, as shown by 
its high number of misclassifications.

Compared to prior studies, the proposed approach achieves comparable or superior performance while relying on substantially fewer sensing channels and features. 
For instance, Trojaniello et al.~\cite{ref30} employed seven kinematic channels, while Pohl et al.~\cite{ref31} used both wrist and ankle IMUs to reach similar 
accuracy levels. Other works, such as those by Zhao et al.~\cite{ref28} and Li et al.~\cite{ref32}, rely on complex pipelines involving inertial navigation or 
signal decomposition, which, although effective, introduce computational and deployment overheads that limit real-world applicability.

In contrast, the proposed feature set is intentionally minimal, interpretable, and biomechanically meaningful. By focusing on extremal angular velocity values 
that directly relate to shank rotation during the gait cycle, the framework preserves discriminative power while reducing sensitivity to sensor noise, feature 
redundancy, and overfitting. These characteristics make the approach particularly suitable for embedded and wearable systems operating under computational, 
energy, and bandwidth constraints.

From an application perspective, the proposed framework is intended for screening-level detection of gait abnormalities in wearable and remote monitoring scenarios, 
where computational efficiency, robustness, and interpretability are prioritized over fine-grained severity stratification. While multi-class severity modeling 
and continuous impairment scoring are important clinical objectives, they require larger, clinically annotated datasets and were therefore considered outside the 
scope of the present study.


Although deep learning approaches such as CNN-TCNs and multilayer perceptrons~\cite{ref33,ref34} have demonstrated strong performance in gait classification, 
their reliance on large datasets and higher computational demands limits their feasibility for real-time rehabilitation and wearable deployment in small-cohort 
clinical studies. The present results demonstrate that, when feature selection is guided by biomechanical relevance and domain knowledge, classical machine 
learning models can achieve strong generalization even under constrained data conditions.

Overall, the results validate that a minimal feature set focused on medically meaningful statistics such as extremal angular velocities of the lower limbs—can 
lead to high predictive performance when paired with appropriately selected classifiers. The combination of simple IMU-derived features, lightweight modeling, 
and subject-level validation highlights a practical pathway toward deployable gait assessment tools for clinical screening and remote rehabilitation support.
